% !TEX root = root.tex
% !TEX recipe = LuaLaTeX+se

\section{Tempus Nativitatis}
\subsection{In Nativitate Domini}
\subsubsection{At Vigil Mass}~

\greannotation{CO. I}
\grecommentary{Cf. Is. 40:5}
\gregorioscore{revelabitur/solesmes.gabc}
%\gregorioscore{revelabitur/english.gabc}
\gresetinitiallines{0}
\grecommentary{Ps. 23}
% 1. 2. 3. 4. 5. 6. 7. 8.
\gregorioscore{revelabitur/verses.gabc}
\gresetinitiallines{1}


\subsubsection{At Mass at Night}
\greannotation{CO. VI}
\grecommentary{Ps. 109:3}
\gregorioscore{in-splendoribus/solesmes.gabc}
%\gregorioscore{in-splendoribus/english.gabc}
\gresetinitiallines{0}
\grecommentary{Ps. 109}
% 1a. 1b. 2. 3. 4. 5. 7.
\gregorioscore{in-splendoribus/verses.gabc}
\gresetinitiallines{1}


\subsubsection{At Mass at Dawn}
\greannotation{CO. IV}
\grecommentary{Zach. 9:9}
\gregorioscore{exsulta-filia/solesmes.gabc}
%\gregorioscore{exsulta-filia/english.gabc}
\gresetinitiallines{0}
\grecommentary{Ps. 33*}
% It just says 33* (differentia g) ?
\gregorioscore{exsulta-filia/verses.gabc}\label{co:exsulta-filia}
\gresetinitiallines{1}


\subsubsection{At Mass at Daytime}
\greannotation{CO. I}
\grecommentary{Ps. 97:3cd}
\gregorioscore{viderunt-omnes/solesmes.gabc}
%\gregorioscore{viderunt-omnes/english.gabc}
\gresetinitiallines{0}
\grecommentary{Ps. 97}
% Ps. 97, I ab. I cd. 2. 3 ab. 4· 5· 6. 7· 8-9 a. 9 be
\gregorioscore{viderunt-omnes/verses.gabc}
\gresetinitiallines{1}


\subsection{Within the Octave of the Nativity}
\emph{All weekdays as on the Nativity at ``Mass at Dawn''
 or ``Mass at Daytime'', except:}
\paragraph{December 29th}
CO. Responsum,~\ref{co:responsum}.

\subsection{Holy Family of Jesus, Mary \& Joseph}
Celebrated on the Sunday within the Octave of the Nativity,
\emph{or} December 30, if the Nativity was on a Sunday.
\paragraph{Year A:}~

\greannotation{CO. VII}
\grecommentary{Mt. 2:20}
\gregorioscore{tolle-puerum/solesmes.gabc}
% %\gregorioscore{tolle-puerum/english.gabc}
% \gresetinitiallines{0}
% \grecommentary{Ps. 92*}
% %Ps. 92*, 1 ab. 1 c-2 a. 3· 4· 5
% %vel ps. 127*, 1. 2. 3· 4· 5· 6
% \gregorioscore{tolle-puerum/verses.gabc}
% \gresetinitiallines{1}

\paragraph{Years B \& C:}~

\greannotation{CO. I}
\grecommentary{Lk. 2:48, 49}
\gregorioscore{fili-quid/solesmes.gabc}
%\gregorioscore{fili-quid/english.gabc}
\gresetinitiallines{0}
\grecommentary{Ps. 26*}
%Ps. 26*, 1 a. 4 abc. 4 de. 5. 8. 13
\gregorioscore{fili-quid/verses.gabc}
\gresetinitiallines{1}


\subsection{Solemnity of Mary, Mother of God}
\emph{N.B.: The antiphon is as on the Nativity, at ``Mass at Dawn'',
but the verses are taken from the Magnificat.}
\greannotation{CO. IV}
\grecommentary{Zach. 9:9}
\gregorioscore{exsulta-filia/solesmes.gabc}
%\gregorioscore{exsulta-filia/english.gabc}
\gresetinitiallines{0}
\grecommentary{Luke 1:46-55}
\gregorioscore{exsulta-filia/magnificat-verses.gabc}
\gresetinitiallines{1}


\subsection{The Epiphany of the Lord}
\greannotation{CO. IV}
\grecommentary{Cf. Mt. 2:2}
\gregorioscore{vidimus-stellam/solesmes.gabc}
%\gregorioscore{vidimus-stellam/english.gabc}
\gresetinitiallines{0}
\grecommentary{Ps. 71*}
%Ps. 71*, 1. 2. 3. 7. 8. 10. 11. 12. 17 ab. 17 cd. 18
\gregorioscore{vidimus-stellam/verses.gabc}
\gresetinitiallines{1}

\paragraph{Weekdays after Epiphany.} As on Epiphany, or on the Nativity
of the Lord at Mass at Dawn.

\paragraph{January 7, where Epiphany is to be celebrated on Sunday, January 8th.}
CO. Dicit Dominus: Implete hydrias,~\ref{co:dicit-dominus-implete}.

\paragraph{January 8, or Tuesday after Epiphany}
CO. Manducaverunt,~\ref{co:Manducaverunt}.

\subsection{The Baptism of the Lord}
Ordinarily celebrated on the Sunday after Epiphany.
When the solemnity of the Epiphany of the Lord is transferred to a
Sunday falling on January 7 or 8, the feast of the Baptism of the Lord is
celebrated on the following Monday.
%\footnote{It is truly
%remarkable the lengths we will go to avoid just one more obligatory
%Mass during the year.}
\greannotation{CO. II}
\grecommentary{Gal. 3:27}
\gregorioscore{omnes-qui-in-christo/solesmes.gabc}
%\gregorioscore{omnes-qui-in-christo/english.gabc}
\gresetinitiallines{0}
\grecommentary{Ps. 28*}
%Ps. 28*, 1. 2. 3. 4. 5. 7-8. 10. 11
\gregorioscore{omnes-qui-in-christo/verses.gabc}
\gresetinitiallines{1}
